% =============================================================================
% Appendix G --- Selected RTL Source Listings
% Every listing is a verbatim quote of the corresponding source file.
% Line numbers match the lines in the repository at the time this report
% was written.
% =============================================================================

\chapter{Selected RTL Source Listings}
\label{app:listings}

This appendix provides verbatim source extracts for the parts of the
RTL that are discussed most frequently in the main chapters. Each
listing is quoted from the exact file path shown in its caption so
that the reader can cross-check any detail.

\section{One-Hot Predicates (\sv{parallel\_interface\_pkg.sv})}
\label{app:listing-onehot}

\begin{lstlisting}[style=svstyle,
    caption={One-hot structural predicates and tail validity, from
             \sv{source/Parallel\_Interface/parallel\_interface\_pkg.sv}.},
    label={lst:onehot}]
function automatic logic pi_is_onehot4(input logic [3:0] oh);
    unique case (oh)
        4'b0001, 4'b0010, 4'b0100, 4'b1000: return 1'b1;
        default: return 1'b0;
    endcase
endfunction

function automatic logic pi_is_onehot8(input logic [7:0] oh);
    unique case (oh)
        8'b00000001, 8'b00000010, 8'b00000100, 8'b00001000,
        8'b00010000, 8'b00100000, 8'b01000000, 8'b10000000: return 1'b1;
        default: return 1'b0;
    endcase
endfunction

// Host ann-tail validity: each field must be strictly one-hot.
function automatic logic ann_tail_is_valid_onehot(input logic [23:0] tail);
    return pi_is_onehot4(tail[23:20]) &&
           pi_is_onehot4(tail[19:16]) &&
           pi_is_onehot8(tail[15:8])  &&
           pi_is_onehot8(tail[7:0]);
endfunction
\end{lstlisting}

\section{Tail-to-Address Decode (\sv{parallel\_interface\_pkg.sv})}
\label{app:listing-addr-decode}

\begin{lstlisting}[style=svstyle,
    caption={Packed-address decode from the 24-bit one-hot tail.},
    label={lst:addr-decode}]
// Decode ann_core_word-style tail -> 16-bit address field
// (parse_ann_address layout)
function automatic logic [15:0] ann_tail_to_parallel_addr(logic [23:0] tail);
    logic [1:0] blk, sb;
    logic [2:0] rid, cid;
    blk = pi_onehot4_to_idx(tail[23:20]);
    sb  = pi_onehot4_to_idx(tail[19:16]);
    cid = pi_onehot8_to_idx(tail[15:8]);
    rid = pi_onehot8_to_idx(tail[7:0]);
    return {6'b0, blk, sb, cid, rid};
endfunction
\end{lstlisting}

\section{Address Pack (Indices \(\rightarrow\) 32-bit Host Word)}
\label{app:listing-addr-pack}

\begin{lstlisting}[style=svstyle,
    caption={Inverse: build a 32-bit host word from a payload byte and
             packed 16-bit parallel address (used by both RTL paths and
             the testbenches).},
    label={lst:addr-pack}]
function automatic logic [31:0] build_host_ann_word(
    input logic [7:0]  data_byte,
    input logic [15:0] parallel_addr
);
    logic [1:0] blk, sb;
    logic [2:0] rid, cid;
    logic [3:0] pe_oh, sa_oh;
    logic [7:0] col_oh, row_oh;
    blk = parallel_addr[9:8];
    sb  = parallel_addr[7:6];
    cid = parallel_addr[5:3];
    rid = parallel_addr[2:0];
    pe_oh  = 4'(1 << blk);
    sa_oh  = 4'(1 << sb);
    col_oh = 8'(1 << cid);
    row_oh = 8'(1 << rid);
    return {data_byte, pe_oh, sa_oh, col_oh, row_oh};
endfunction
\end{lstlisting}

\section{Pulse-Train Helpers (\sv{controller\_pkg.sv})}
\label{app:listing-pulse-helpers}

\begin{lstlisting}[style=svstyle,
    caption={Pulse-train total/active helpers, from
             \sv{source/Controller/controller\_pkg.sv}.},
    label={lst:pulse-helpers}]
// Burst-train helpers (cycle_idx is 0 .. pulse_train_total-1)
function automatic int pulse_train_total(input int T, input int N, input int G);
    int t_safe;
    t_safe = (T < 1) ? 1 : T;
    if (N < 1)
        return 0;
    return N * t_safe + (N - 1) * G;
endfunction

function automatic logic pulse_train_active(
    input int cycle_idx,
    input int T,
    input int N,
    input int G
);
    int t_safe, chunk, r, base;
    t_safe = (T < 1) ? 1 : T;
    if (N < 1 || cycle_idx < 0)
        return 1'b0;
    chunk = t_safe + G;
    for (r = 0; r < N; r++) begin
        base = r * chunk;
        if (cycle_idx >= base && cycle_idx < base + t_safe)
            return 1'b1;
    end
    return 1'b0;
endfunction
\end{lstlisting}

\begin{lstlisting}[style=svstyle,
    caption={LUT-based macro-repeat helpers.},
    label={lst:pulse-lut}]
// R back-to-back copies of one macro train (length Mmacro),
// separated by G HIZ cycles (no gap after the last copy).
// Used for LUT first PROG so repeats are visible when N=1.
function automatic int pulse_lut_macro_repeat_total(
    input int R, input int Mmacro, input int G
);
    int Rc;
    Rc = (R < 1) ? 1 : R;
    if (Mmacro < 1)
        return 0;
    return Rc * Mmacro + (Rc - 1) * G;
endfunction

function automatic logic pulse_lut_macro_repeat_active(
    input int cyc,
    input int R,
    input int Mmacro,
    input int T,
    input int N,
    input int G
);
    int Rc;
    int cum;
    int b;
    Rc = (R < 1) ? 1 : R;
    if (Mmacro < 1 || cyc < 0)
        return 1'b0;
    cum = 0;
    for (b = 0; b < Rc; b++) begin
        if (cyc >= cum && cyc < cum + Mmacro)
            return pulse_train_active(cyc - cum, T, N, G);
        cum = cum + Mmacro + G;
    end
    return 1'b0;
endfunction
\end{lstlisting}

\section{\texttt{VERIFY\_CHECK} Case Body (\sv{controller.sv})}
\label{app:listing-verify-check}

\begin{lstlisting}[style=svstyle,
    caption={\sv{VERIFY\_CHECK} branch of the combinational FSM block,
             from \sv{source/Controller/controller.sv}.},
    label={lst:verify-check}]
VERIFY_CHECK: begin
    // Compare read weight with expected weight
    if (weight_read_data == expected_weight) begin
        verify_next = VERIFY_DONE;   // Match: continue to next weight
    end else if (weight_read_data < expected_weight) begin
        // Under-programmed: re-apply PROG pulses (no ERASE)
        if (prog_retry_cnt < controller_pkg::MAX_PROG_RETRIES) begin
            next_state = S_PROGRAM;
            next_prog_state = PROG_HIZ;
        end else begin
            // Max re-prog attempts exceeded, fall back to ERASE
            verify_failure_starts_erase = 1'b1;
            next_state = S_ERASE;
        end
        verify_next = VERIFY_IDLE;
    end else begin
        // Over-programmed: ERASE then retry
        verify_failure_starts_erase = 1'b1;
        next_state = S_ERASE;
        verify_next = VERIFY_IDLE;
    end
end
\end{lstlisting}

\section{Bit-Serial Output Block (\sv{input\_buffer.sv})}
\label{app:listing-bit-serial}

\begin{lstlisting}[style=svstyle,
    caption={Bit-serial lane generation in \modname{input\_buffer}.
             From \sv{source/Input\_Buffer/input\_buffer.sv}.},
    label={lst:bit-serial}]
// Read Path: D0-D7 = 1 bit per channel per cycle (bit-serial, LSB-first)
// For computation (CTRL_COMPUTE), controller drives bit_sel 0..7 each cycle.
// D0 = buffer_reg[addr][bit_sel], D1 = buffer_reg[addr+1][bit_sel], ...
`comb(
    D0 = 1'b0;
    D1 = 1'b0;
    D2 = 1'b0;
    D3 = 1'b0;
    D4 = 1'b0;
    D5 = 1'b0;
    D6 = 1'b0;
    D7 = 1'b0;

    if (read_en) begin
        if (is_valid_addr(addr))
            D0 = buffer_reg[addr][bit_sel];
        if ((addr + 1) < BUFFER_SIZE)
            D1 = buffer_reg[addr + 1][bit_sel];
        if ((addr + 2) < BUFFER_SIZE)
            D2 = buffer_reg[addr + 2][bit_sel];
        if ((addr + 3) < BUFFER_SIZE)
            D3 = buffer_reg[addr + 3][bit_sel];
        if ((addr + 4) < BUFFER_SIZE)
            D4 = buffer_reg[addr + 4][bit_sel];
        if ((addr + 5) < BUFFER_SIZE)
            D5 = buffer_reg[addr + 5][bit_sel];
        if ((addr + 6) < BUFFER_SIZE)
            D6 = buffer_reg[addr + 6][bit_sel];
        if ((addr + 7) < BUFFER_SIZE)
            D7 = buffer_reg[addr + 7][bit_sel];
    end
)
\end{lstlisting}

\section{Simulation Flow Script Snippet (\sv{run\_sim.py})}
\label{app:listing-runsim}

\begin{lstlisting}[style=plainstyle,language={},
    caption={Representative commands from
             \sv{scripts/run\_sim.py}. The script wraps \sv{vlog}
             (compilation) and \sv{vsim} (simulation), and supports
             \sv{list}, \sv{compile}, \sv{sim}, \sv{run\_all}, and
             \sv{clean} subcommands.},
    label={lst:runsim}]
# compile verification bundle
python scripts/run_sim.py compile -m Controller -t verif
python scripts/run_sim.py compile -m Input_Buffer -t verif
python scripts/run_sim.py compile -m Parallel_Interface -t verif

# run a single testbench
python scripts/run_sim.py sim -m Controller \
    -tb controller_prog_verify_lut_tb

# run every testbench for a module
python scripts/run_sim.py run_all -m Controller

# open a wave wrapper in the GUI
python scripts/run_sim.py sim -m Controller \
    -tb controller_prog_verify_lut_tb_waves_tb \
    --do-file verif/Controller/do/waves/controller_prog_verify_lut_tb_waves.do
\end{lstlisting}
